\documentclass[main.tex]{subfiles}


\begin{document}

%from pagenumber to up
% \phantomsection 
% \hypertarget{toc}{}

 

 

% номер секции вручную
\setcounter{section}{30
}
\addtocounter{section}{-1} 

\section{Сложение магнитных полей}


Магнитное поле, создаваемое несколькими магнитами, можно рассматривать как наложение полей, создаваемых каждым магнитом в отдельности.
%
\begin{law}
\textbf{Принцип суперпозиции магнитных полей.} Если магниты $\text{М}_1, \text{М}_2,\ldots$ по отдельности создают в данной точке поля $\vec B_1,\vec B_2,\ldots,$ то вместе они создают в данной точке поле
\begin{equation}
    \label{34Bprinc_e1}
    \vec B=\vec B_1+\vec B_2+\ldots
\end{equation}
\end{law}
%
Этот принцип можно проиллюстрировать для случая двух магнитов (рис.\,\ref{fig:p142}).

    \begin{figure}[H]
    \centering

    \begin{tikzpicture}[font=\small,            line cap=round,                             >={Stealth[inset=0pt,round,length=3mm, angle'=20]},vec/.style={>={Stealth[inset=0pt,round,length=3.5mm, angle'=20]}}]


\def\lm{1} 


%magnet
\fill[red, semitransparent] (0,0) rectangle++ (\lm,{\lm/2});
\fill[NavyBlue, semitransparent] (\lm,0) rectangle++ (\lm,{\lm/2});

\node[right] at (0,.25) {$S$}; 
\node[left] at (2,.25) {$N$}; 

\draw (0,0) rectangle++ ({2*\lm},{\lm/2});

\node[below] at (\lm,0) {$\text{М}_1$}; 

%right magnet
\begin{scope}[xshift={2.5*\lm cm},yshift=-{\lm/2*1cm}]
\fill[NavyBlue, semitransparent] (0,0) rectangle++ ({\lm/2},-\lm);
\fill[red, semitransparent] (0,-\lm) coordinate (mb) rectangle++ ({\lm/2},-\lm);

\node[above] at (.25,-2) {$S$}; 
\node[below] at (.25,.) {$N$}; 

\draw (0,0) rectangle++ ({\lm/2},{-2*\lm});

\node[left] at (mb) {$\text{М}_2$}; 

\end{scope}


% point
\draw[gray,->] ({2.75*\lm},{\lm/4}) coordinate (p) --++ (1,0) node[black,below] {$\vec B_1$};
\draw[gray,->] (p) --++ (0,1) node[black,left] {$\vec B_2$};

\draw[dashed,gray] (p) ++ (0,1) --++ (1,0) coordinate (b) --++ (0,-1);

\draw[thick,vec,gray,->] (p) -- (b) node[black,right] {$\vec B$};
    
\node at (p) [circle,fill=,inner sep=0pt,minimum size=1mm,label={below left,inner xsep=0:$A$}] {};
    
    
    \end{tikzpicture}
\caption{Принцип суперпозиции магнитных полей}
\label{fig:p142}
\end{figure}

Магнит~$\text{М}_1$ создает в точке~$A$ поле~$\vec B_1$, а магнит~$\text{М}_2$ в этой же точке создает поле~$\vec B_2$. Согласно вышеуказанному принципу вместе они создают в точке~$A$ поле~$\vec B=\vec B_1+\vec B_2$ (рис.\,\ref{fig:p142}). 

\emph{Индукции магнитных полей в общем случае складываются векторно.}

\medskip

\noindent\textbf{Задача.} Два одинаковых магнита~$\text{М}_1,\text{М}_2$ расположены так, что их оси лежат на одной прямой~$OO'$; магниты обращены друг к другу одноименными полюсами (рис.\,\ref{fig:p143}). Чему равен модуль результирующего вектора магнитной индукции в точке~$A$, лежащей на прямой~$OO'$ и равноудаленной от обоих магнитов? 
% Как будет вести себя магнитная стрелка помещенная в точку~$A$?  

    \begin{figure}[H]
    \centering

    \begin{tikzpicture}[font=\small,            line cap=round,                             >={Stealth[inset=0pt,round,length=3mm, angle'=20]},vec/.style={>={Stealth[inset=0pt,round,length=3.5mm, angle'=20]}}]


\def\lm{1} 


%magnet
\fill[red, semitransparent] (0,0) rectangle++ (\lm,{\lm/2});
\fill[NavyBlue, semitransparent] (\lm,0) rectangle++ (\lm,{\lm/2});

\node[right] at (0,.25) {$S$}; 
\node[left] at (2,.25) {$N$}; 

\draw (0,0) rectangle++ ({2*\lm},{\lm/2});

\node[below] at (\lm,0) {$\text{М}_1$}; 

%right magnet
\begin{scope}[xshift={4*\lm cm}]
\fill[NavyBlue, semitransparent] (0,0) rectangle++ (\lm,{\lm/2});
\fill[red, semitransparent] (\lm,0) rectangle++ (\lm,{\lm/2});

\node[right] at (0,.25) {$N$}; 
\node[left] at (2,.25) {$S$}; 

\draw (0,0) rectangle++ ({2*\lm},{\lm/2});

\node[below] at (\lm,0) {$\text{М}_2$}; 

\end{scope}


%dashed
\draw[gray,dashed] (0,{\lm/4}) --++ (-1,0) node[black,left] {$O$};
\draw[gray,dashed] ({6*\lm},{\lm/4}) --++ (1,0) node[black,right] {$O'$};
\draw[gray,dashed] ({2*\lm},{\lm/4}) --++ ({2*\lm},0);


% point
\node at ({3*\lm},{.25*\lm}) [circle,fill=,inner sep=0pt,minimum size=1mm,label={above,inner xsep=0:$A$}] {};

    
    \end{tikzpicture}
\caption{К задаче}
\label{fig:p143}
\end{figure}

\noindent\emph{Решение.} Так как точка~$A$ одинаково удалена от одноименных полюсов магнитов~$\text{М}_1$ и~$\text{М}_2$ (или густота собственных магнитных линий каждого магнита в этой точке одинакова), то соответствующие магнитные индукции~$\vec B_1$ и~$\vec B_2$, создаваемые этими магнитами в данной точке, одинаковы по величине (рис.\,\ref{fig:p144}).

    \begin{figure}[H]
    \centering

    \begin{tikzpicture}[font=\small,            line cap=round,                             >={Stealth[inset=0pt,round,length=3mm, angle'=20]},vec/.style={>={Stealth[inset=0pt,round,length=3.5mm, angle'=20]}}]


\def\lm{1} 


%dashed
\draw[gray,dashed] ({4*\lm},{\lm/4}) --++ ({-4*\lm},0) node[black,left] {$O$} ++ ({4*\lm},0) --++ ({4*\lm},0) node[black,right] {$O'$};


%vec
\draw[thick,vec,gray,->] ({4*\lm},{.25*\lm}) --++ (1,0) node[black,above] {$\vec B_1$};
\draw[thick,vec,gray,->] ({4*\lm},{.25*\lm}) --++ (-1,0) node[black,above] {$\vec B_2$};


% point
\node at ({4*\lm},{.25*\lm}) [circle,fill=,inner sep=0pt,minimum size=1mm,label={above,inner xsep=0:$A$}] {};

    
    \end{tikzpicture}
\caption{К задаче}
\label{fig:p144}
\end{figure}

Так как векторы~$\vec B_1$ и~$\vec B_2$ равны по модулю и противоположны по направлению, то результирующий вектор~$\vec B$, вычисляемый по формуле~\eqref{34Bprinc_e1}, равен нулю: $\vec B=\vec B_1+\vec B_2=0$. Соответственно, модуль вектора~$\vec B$ равен нулю.  














 
\end{document}
